\chapter{Grundlæggende aritmetik}
\label{chap:ch0}
En solid forståelse af de grundlæggende regneoperationer er fundamentet for mere avancerede matematiske emner. Aritmetik handler om at behandle tal ved hjælp af grundlæggende operationer som addition, subtraktion, multiplikation og division. Disse operationer er ikke kun vigtige i hverdagens beregninger. De er også byggestenene i mere komplekse matematiske metoder.

I dette kapitel gennemgår vi de centrale teknikker inden for algebraisk regning: faktorisering og regnearternes hierarki, regning med brøker, potenser og rødder, brug af formler og indsættelse samt omskrivning af formler. Når du behersker disse færdigheder, er du klar til at studere funktioner og de emner, der følger.

\section{Faktorisering og regnearternes hierarki}
Husk, at når vi beregner et produkt som

\[
a(b + c) = ab + ac
\]
ganger vi faktoren \(a\) ind på hvert led i parentesen. Dette kalder vi \textbf{den distributive lov}. Det er dog ofte en fordel eller ligefrem nødvendigt at udføre den omvendte operation, som kaldes faktorisering. Ved faktorisering skriver vi summen \(ab + ac\) på faktoriseret form \(a(b + c)\).

Matematisk betragtes udtrykket \(a(b + c)\) som mere reduceret eller »bedre« end \(ab + ac\). For at forstå hvorfor må vi skelne mellem \textbf{led} og \textbf{faktorer}. Led er adskilt af addition eller subtraktion, mens faktorer er adskilt af multiplikation eller division. For eksempel består udtrykket

\[ 
8 - 5 + 3 
\]
af tre led (8, -5 og 3), mens udtrykket

\[ 
2 \times 5 \times 3 
\]
består af tre faktorer (2, 5 og 3). Se på udtrykket

\[ 
5 + 2 + 7 \times 8 
\]
som består af tre led (5, 2 og \(7 \times 8\)). Bemærk, at et af leddene selv indeholder to faktorer (7 og 8). På samme måde består udtrykket

\[ 
5(2 + 4)(-9) 
\]
af tre faktorer (5, \(2 + 4\) og -9), hvor faktoren \(2 + 4\) indeholder to led.

Fordelen ved at skrive matematiske udtryk på faktoriseret form er, at de så kan reduceres mere effektivt. Det bliver tydeligere i løbet af kapitlet, især når vi arbejder med brøker og ligninger. Her er nogle eksempler på faktorisering:

\newpage

\begin{example} Faktorisering

\begin{align*}
ab - ac &= a(b - c) \\
-ab - ac &= a(-b - c) \\
-ab - ac &= -a(b + c) \\
ab - ac + aa - aa &= a(b - c + a - a) \\
abc - ab + aba &= ab(c - 1 + a) \\
ab - abc - a &= a(b - bc - 1)
\end{align*}

\end{example}

Når vi udregner matematiske udtryk, er det afgørende at følge en bestemt rækkefølge for operationerne for at få korrekte resultater. Operationerne udføres i denne rækkefølge:

\begin{enumerate}
    \item \textbf{Parenteser}: Udfør først alle operationer inde i parenteser.
    \item \textbf{Potenser og rødder}: Udregn derefter potenser og rødder.
    \item \textbf{Multiplikation og division}: Udfør så multiplikation og division fra venstre mod højre i den rækkefølge, de optræder.
    \item \textbf{Addition og subtraktion}: Udfør til sidst addition og subtraktion fra venstre mod højre i den rækkefølge, de optræder.
\end{enumerate}

Lad os se på et eksempel for hver operation for at illustrere regnearternes hierarki:

\begin{example} Parenteser \newline
Udregn udtrykket: \((2 + 3) \times 4\)

\[
(2 + 3) \times 4 = 5 \times 4 = 20
\]
\end{example}

\begin{example} Potenser og rødder \newline
Udregn udtrykket: \(3^2 + \sqrt{16}\)

\[
3^2 + \sqrt{16} = 9 + 4 = 13
\]
\end{example}

\begin{example} Multiplikation og division \newline
Udregn udtrykket: \(6 \div 2 \times 3\). Ifølge det almindelige regnearternes hierarki udfører vi division og multiplikation fra venstre mod højre:

\[
6 \div 2 \times 3 = (6 \div 2) \times 3 = 3 \times 3 = 9
\]

\end{example}

\begin{example} Addition og subtraktion \newline
Udregn udtrykket: \(8 - 3 + 2\)

\[
8 - 3 + 2 = 5 + 2 = 7
\]
\end{example}

\section{Brøker}
Per definition består en brøk altid af (mindst) to faktorer. Den første faktor kalder vi \textbf{tælleren}. Det er brøkens »øverste del«. Den nederste del kalder vi \textbf{nævneren}. Den måske vigtigste regel for regning med brøker er, at to brøker kun kan lægges sammen eller trækkes fra hinanden, hvis de har samme nævner. Desuden må nævneren aldrig være lig med 0.

\begin{example}

\[
\frac{1}{x^2 - 2}
\]    
\end{example}

Her skal vi sikre, at \(x^2 - 2 \neq 0\). Det betyder, at vi kun må bruge værdier, der er forskellige fra \(\pm\sqrt{2}\).

\begin{proposition}[Regneregler for brøker]
For $a,b,c,m \in \mathbb{R}$, hvor $a,b,c,m \neq 0$ når det er nødvendigt, gælder følgende identiteter:
\begingroup
\setlength{\jot}{8pt} % increase space between align rows (default ~3pt)
\begin{align*}
(1) & \quad \frac{a}{b} \times m = \frac{am}{b} \\
(2) & \quad \frac{a}{b} \div m = \frac{a}{bm} \\
(3) & \quad m \div \frac{a}{b} = \frac{mb}{a} \\
(4) & \quad \frac{a}{b} \times \frac{c}{a} = \frac{c}{b} \\
(5) & \quad \frac{a}{b} \div \frac{c}{a} = \frac{a^2}{bc} \\
(6) & \quad \frac{a}{b} = \frac{ac}{bc} \\
(7) & \quad \frac{a}{b} + \frac{c}{a} = \frac{a^2 + bc}{ab}
\end{align*}
\endgroup
\end{proposition}

For at forlænge eller forkorte en brøk skal vi gange eller dividere med det samme tal i både tæller og nævner:

\begin{example} Forlængelse og forkortning af brøker

\[
\frac{72}{144} = \frac{72 \div 12}{144 \div 12} = \frac{6}{12} = \frac{6 \div 6}{12 \div 6} = \frac{1}{2}
\]

\[
\frac{2}{3} = \frac{2 \times 6}{3 \times 6} = \frac{12}{18}
\]

\[
\frac{2x}{2xx} = \frac{2}{2x}
\]

\[
\frac{3a}{6a + 3b} = \frac{3a}{3(2a + b)} = \frac{a}{2a + b}
\]    
\end{example}

Når vi faktoriserer et udtryk, skal alle led divideres eller ganges på samme måde. Det betyder, at følgende udtryk ikke kan reduceres yderligere, selv om det kan være fristende:

\[
\frac{a}{2a + b} \neq \frac{1}{2 + b}
\]

For at faktorisere \(\frac{a}{2a + b}\) korrekt skal du dividere \(a\) op i alle led i nævneren:

\[
\frac{a}{2a + b} = \frac{1}{2 + \frac{b}{a}}
\]

Som vist ovenfor er multiplikation af brøker ligetil: Man ganger tællerne med hinanden og nævnerne med hinanden.

\begin{example} Multiplikation af brøker \newline
Se på brøkerne \(\frac{a}{b}\) og \(\frac{c}{d}\). Deres produkt er:

\[
\frac{a}{b} \times \frac{c}{d} = \frac{a \times c}{b \times d}
\]

Hvis for eksempel \(a = 2\), \(b = 3\), \(c = 4\) og \(d = 5\), så er:

\[
\frac{2}{3} \times \frac{4}{5} = \frac{2 \times 4}{3 \times 5} = \frac{8}{15}
\]
  
\end{example} 

\begin{example} Division af brøker

\[
\frac{6}{7} \div \frac{4}{21} = \frac{6}{7} \times \frac{21}{4} = \frac{126}{28} = \frac{63}{14} = \frac{9}{2}
\]

\[
\frac{1}{2} \div 3 = \frac{1}{2} \div \frac{3}{1} = \frac{1}{2} \times \frac{1}{3} = \frac{1}{6}
\]

\[
\frac{2}{3} \div \frac{8}{9} = \frac{2}{3} \times \frac{9}{8} = \frac{18}{24} = \frac{3}{4}
\]

\[
\frac{8}{9} \div 16 = \frac{8}{9} \times \frac{1}{16} = \frac{8}{144} = \frac{4}{72} = \frac{1}{18}
\]
  
\end{example}

Addition og subtraktion af brøker volder ofte flere problemer end multiplikation og division. Nøglen er at finde en fællesnævner for brøkerne og huske reglen ovenfor om forlængelse af brøker.

\begin{example} Addition og subtraktion af brøker

 \[
\frac{1}{5} + \frac{2}{5} = \frac{1 + 2}{5} = \frac{3}{5}
\]

\[
\frac{1}{4} + \frac{2}{3} = \frac{1 \times 3}{4 \times 3} + \frac{2 \times 4}{3 \times 4} = \frac{3}{12} + \frac{8}{12} = \frac{3 + 8}{12} = \frac{11}{12}
\]

\[
\frac{7}{12} - \frac{5}{8} = \frac{7 \times 8}{12 \times 8} - \frac{5 \times 12}{8 \times 12} = \frac{56}{96} - \frac{60}{96} = \frac{56 - 60}{96} = \frac{-4}{96} = \frac{-1}{24}
\]

\end{example}

Bemærk: Gør aldrig sådan:

\[
\frac{a}{b} + \frac{c}{d} = \frac{a + b}{b + d}
\]   

\begin{example}

\[
\frac{x}{x - 1} \times \frac{2}{x(x + 4)} = \frac{2x}{(x - 1)x(x + 4)} = \frac{2}{(x - 1)(x + 4)}
\]

\[
\frac{2}{x - 1} \div \frac{x}{x - 1} = \frac{2}{x - 1} \times \frac{x - 1}{x} = \frac{2(x - 1)}{(x - 1)x} = \frac{2}{x}
\]

\[
\frac{x + 1}{x^2 + 2} + \frac{x - 6}{x^2 + 2} = \frac{x + 1 + x - 6}{x^2 + 2} = \frac{2x - 5}{x^2 + 2}
\]    
\end{example}

Husk, at når der står et minustegn foran en brøk, skal alle fortegn i tælleren ændres. Det svarer til, at man ændrer alle fortegn inde i en parentes, når der står et minustegn foran den.

Se på udtrykket:

\[
-\frac{a - b}{c}
\]

For at håndtere minustegnet korrekt ændrer vi alle fortegn i tælleren:

\[
-\frac{a - b}{c} = \frac{-a + b}{c}
\]

Hvis for eksempel \(a = 5\) og \(b = 3\), så er:

\[
-\frac{5 - 3}{c} = \frac{-5 + 3}{c} = \frac{-2}{c}
\]

\begin{example}

\[
\frac{x + 1}{x^2 + 2} - \frac{x - 6}{x^2 + 2} = \frac{x + 1 - x + 6}{x^2 + 2} = \frac{7}{x^2 + 2}
\]    
\end{example}

\section{Potenser, rødder og irrationelle rødder}
En \textbf{eksponent} er en kort måde at skrive gentagen multiplikation af det samme tal på:

\begin{example} Potensopløftning

\[
4 \times 4 \times 4 \times 4 \times 4 = 4^5
\]

\[
x \times x \times x \times x \times x = x^5
\]  
\end{example}

\textbf{Rødder} er tæt knyttet til potensopløftning. Over de reelle tal er kvadratroden \(\sqrt{x}\) defineret som det ikke-negative tal, hvis kvadrat er \(x\). Den er derfor kun defineret, når \(x\geq 0\). Hvis man kvadrerer et reelt tal og derefter tager kvadratroden, får man følgelig tallets absolutte værdi: \(\sqrt{x^2}=|x|\).

\begin{example} Uddragning af rod

\[
\sqrt{a} \times \sqrt{a} = (\sqrt{a})^2 = a
\]

\[
\sqrt{a} = b \iff a = b^2 \text{ og } b\geq 0
\]    
\end{example}



En \textbf{irrationel rod} (på engelsk \emph{surd}) er en rod, der både er reel og irrational. Eksempler er \(\sqrt{2}, \sqrt{3}, \sqrt{5}\) og \(\sqrt{6}\).

Tal kan opløftes i andre potenser end 2, for eksempel i tredje potens eller i fjerde potens, 100. potens og så videre. Tilsvarende kan man tage kubikroden af et tal, fjerderoden, den 100. rod og så videre. For at angive en anden rod end kvadratroden bruger man det samme rodtegn, men indsætter et tal, kaldet rodeksponenten, i rodtegnet. Tallet placeres typisk i den del af tegnet, der ligner et »flueben«.

\begin{example} Rodeksponent og argument

\[
4^3 = 64 \iff \sqrt[3]{64} = 4
\]
   
\end{example}

I dette eksempel er »3« i rodtegnet rodens \textbf{rodeksponent}. »64« kaldes argumentet for \textbf{roden}, også kendt som \textbf{radikanden}. Da kvadratrødder er den mest almindelige type rødder, udelader man normalt rodeksponenten for kvadratrødder. Selv om »\(\sqrt[2]{2}\)« teknisk set er korrekt, bruges det sjældent i praksis.
\begin{proposition}{Regneregler for rødder}
    \begin{align*}
    &(1) \quad \sqrt[n]{x y}=\sqrt[n]{x} \times \sqrt[n]{y}   &   &\textnormal{hvor } x,y \geq 0\\
    &(2) \quad \sqrt[n]{\frac{x}{y}} = \frac{\sqrt[n]{x}}{\sqrt[n]{y}}    &   &\textnormal{hvor } x \geq 0 \textnormal{ og } y > 0 \\
    &(3) \quad \sqrt{x^2} = |x|   &   &\textnormal{hvor } x \in \mathbb{R} \\
    &(4) \quad (\sqrt[n]{x})^n = x    &   &\textnormal{hvor } x\geq 0, \textnormal{ eller hvor } x<0 \textnormal{ og } n \textnormal{ er ulige}\\
    &(5) \quad \sqrt[n]{-x}=-\sqrt[n]{x}  &   &\textnormal{hvor } x \geq 0 \textnormal{ og } n \in \mathbb{N} \textnormal{ er ulige }
    \end{align*}
\end{proposition}

At opløfte et tal i en \textbf{potens}, også kaldet \textbf{potensopløftning}, er en grundlæggende matematisk operation. Som vi så ovenfor, ganger man et tal med sig selv et bestemt antal gange. \textbf{Grundtallet} er det tal, der ganges, og \textbf{eksponenten} angiver, hvor mange gange grundtallet indgår som faktor. For eksempel betyder \(a^n\), at grundtallet \(a\) ganges med sig selv \(n\) gange. Potensopløftning er et stærkt værktøj i matematikken, og nogle få vigtige regler styrer, hvordan det bruges.

\begin{proposition}{Potensregneregler for heltallige eksponenter}
\label{prop:DefPow}
Lad $n,m\in\mathbb{Z}$. Følgende identiteter gælder, når alle de indgående udtryk er definerede. Især skal et grundtal med negativ eksponent være forskelligt fra nul.
\[
\begin{aligned}
&\text{(1)}\quad x^{n}\cdot x^{m}=x^{\,n+m},\\[2pt]
&\text{(2)}\quad \dfrac{x^{n}}{x^{m}}=x^{\,n-m}\quad \text{med } x\neq 0,\\[2pt]
&\text{(3)}\quad x^{n}\cdot y^{n}=(xy)^{n},\\[2pt]
&\text{(4)}\quad \dfrac{x^{n}}{y^{n}}=\left(\dfrac{x}{y}\right)^{n}\quad \text{med } y\neq 0,\\[2pt]
&\text{(5)}\quad \bigl(x^{n}\bigr)^{m}=x^{\,nm},\\[2pt]
&\text{(6)}\quad x^{1}=x.
\end{aligned}
\]
\end{proposition}


Nogle af disse regler gør det muligt at udvide potensbegrebet, så eksponenten kan være et vilkårligt helt tal. Hvis man sætter \(n = m\) i regel (2), får man:

\[
\frac{x^n}{x^n} = x^{n-n} = x^0
\]

For \(x\neq 0\) er brøken \(\dfrac{x^n}{x^n}\) lig med 1, så vi får

\[
x^0 = 1
\]

Dermed er potensbegrebet udvidet til at omfatte \(n \in \mathbb{N} \cup \{0\}\). Hvis man nu igen sætter \(n = 0\) i regel (2), får man:

\[
\frac{x^0}{x^m} = x^{0-m} = x^{-m}
\]

Men ifølge den forrige udregning er \(x^0 = 1\). Derfor får man

\begin{equation*}
 \frac{1}{x^m} = x^{-m}   
\end{equation*}

Som vist i \autoref{prop:DefPow} giver disse regler en fuldstændig ramme for at regne med udtryk med vilkårlige heltallige eksponenter. Som en konsekvens kan følgende yderligere regler udledes:

\begin{proposition}{Flere potensregneregler for heltallige eksponenter}
Lad $n,m\in\mathbb{Z}$. Så gælder følgende (med $x,y\in\mathbb{R}$ og forskellige fra nul, hvor det er angivet):
\[
\begin{aligned}
&\text{(7)}\quad x^{0}=1   &\qquad&   x \neq 0  \\
&\text{(8)}\quad \frac{1}{x^{m}}=x^{-m} &\qquad& x \neq 0 \\
\end{aligned}
\]
\end{proposition}

Lad os illustrere disse regler med et par eksempler.

\begin{example} Reducér følgende udtryk

    \begin{equation*}
\left(3 x y^6\right)^3=3^3 \cdot x^3 \cdot\left(y^6\right)^3=27 x^3 y^{18}
\end{equation*}
\end{example}

\begin{example} Reducér følgende udtryk

\begin{equation*}
\frac{a^{-4} b^3}{a^7 b^{-5}}=\frac{a^{-4}}{a^7} \cdot \frac{b^3}{b^{-5}}=a^{-4-7} \cdot b^{3-(-5)}=a^{-11} \cdot b^8=\frac{b^8}{a^{11}}
\end{equation*}
\end{example}

For positive tal \(x\) kan potensbegrebet udvides yderligere, så det også gælder, når eksponenten er et rationalt tal. Ethvert rationalt tal \(r \in \mathbb{Q}\) kan skrives som \(r = \dfrac{m}{n}\), hvor \(m \in \mathbb{Z}\) og \(n \in \mathbb{N}\). For \(x > 0\) definerer vi nu

\begin{equation*}
y=x^r=x^{\frac{m}{n}}
\end{equation*}

Heraf får man (blandt andet ved hjælp af regel (5)):

\[
y^n = \left(x^{\frac{m}{n}}\right)^n = x^{\frac{m}{n} \cdot n} = x^m
\]

Ved at bruge rodbegrebet får vi til sidst

\begin{equation*}
 y^n=x^m \quad \iff \quad y=\sqrt[n]{x^m}   
\end{equation*}

Bemærk, at fordi \(x > 0\), er også \(y > 0\). Vi er nu klar til at formulere \textbf{\textit{det udvidede potensbegreb}}.

\begin{definition}{Det udvidede potensbegreb}
$\textnormal{Lad } m \in \mathbb{Z} \textnormal{ og lad }  n \in \mathbb{N}\textnormal{, så } \dfrac{m}{n} \in \mathbb{Q}  \textnormal{. Så gælder følgende:}$

\[
x^{\frac{m}{n}} = \sqrt[n]{x^m} \hspace{1cm} x > 0
\]

\textnormal{Mere specifikt gælder}

\[
x^{\frac{1}{n}} = \sqrt[n]{x} \hspace{1.3cm} x > 0
\]

\end{definition}

Nævneren i en rational eksponent svarer til rodens rodeksponent, mens tælleren forbliver eksponent for grundtallet. Omvendt kan en rods rodeksponent omdannes til nævneren i en eksponent i et tilsvarende potensudtryk. Denne egenskab gør det muligt at omskrive ethvert rodudtryk til potensform, hvilket er et stærkt værktøj til reduktion.

\begin{example}

\[
\sqrt[5]{x^3} = x^{\frac{3}{5}} \quad  \textnormal{mod} \quad \sqrt[3]{x^5} = x^{\frac{5}{3}}
\]

\[
\frac{1}{\sqrt[7]{x^3}} = x^{-\frac{3}{7}}
\]

\[
\frac{1}{\sqrt[3]{x^2}} = \left(x^2\right)^{-\frac{1}{3}} = x^{-\frac{2}{3}}
\]
\end{example}

Denne egenskab kan også vendes om: Enhver rational eksponent kan skrives om til et rodudtryk ved at bruge nævneren som rodens rodeksponent. Når vi kan skifte mellem potensform og rodform, kan vi udregne udtryk, som tidligere var svære at håndtere, ved at omskrive dem til rødder.

\begin{example}
\[
27^{-\frac{4}{3}} = \frac{1}{\sqrt[3]{27^4}} = \frac{1}{\left(\sqrt[3]{27}\right)^4} = \frac{1}{3^4} = \frac{1}{81}
\]    
\end{example}

En af de største fordele ved at omskrive et rodudtryk til potensform er, at vi så kan bruge alle potensreglerne til at reducere udtrykket. De følgende eksempler viser, hvordan forskellige regler kan bruges til at reducere udtryk med rationale eksponenter.

\begin{example}

\[
a^{\frac{2}{3}}b^{\frac{1}{2}}a^{\frac{1}{6}}b^{\frac{1}{5}} = a^{\frac{2}{3}+\frac{1}{6}}b^{\frac{1}{2}+\frac{1}{5}} = a^{\frac{4}{6}+\frac{1}{6}}b^{\frac{5}{10}+\frac{2}{10}} = a^{\frac{5}{6}}b^{\frac{7}{10}}
\]

\[
\left( x^{\frac{1}{3}} x^{\frac{2}{5}} \right)^{\frac{3}{4}} = x^{\frac{1}{3} \times \frac{3}{4}} x^{\frac{2}{5} \times \frac{3}{4}} = x^{\frac{3}{12}} x^{\frac{6}{20}} = x^{\frac{1}{4}} x^{\frac{3}{10}} = x^{\frac{11}{20}}
\]

\[
\frac{x^{\frac{4}{2}}x^{\frac{4}{6}}x^{\frac{1}{2}}x^{\frac{5}{6}}}{x^{\frac{7}{2}}x^0}
=x^{\frac{4}{2}+\frac{4}{6}+\frac{1}{2}+\frac{5}{6}-\frac{7}{2}}
=x^{\frac{1}{2}}
\]

\[
\left(25x^{\frac{1}{3}}x^{\frac{2}{5}}\right)^{-\frac{1}{2}}
=\left(25x^{\frac{11}{15}}\right)^{-\frac{1}{2}}
=\frac{1}{5x^{\frac{11}{30}}}
\]
\end{example}

Det er vigtigt at huske, at når vi reducerer udtryk med rationale eksponenter, bruger vi de samme potensregler som for heltallige eksponenter. Den eneste forskel er, at vi også skal overholde reglerne for brøker.

\section{Brug af formler og indsættelse}
I ingeniørfagene er fysiske størrelser ofte knyttet til hinanden gennem formler. Formlerne består af variable og konstanter, som repræsenterer de pågældende fysiske størrelser. For at udregne en formel skal man indsætte talværdier for variablene.

Ohms lov er for eksempel en formel, der forbinder spændingen \(v\) over en modstand med resistansen \(R\) og strømmen \(i\), der løber gennem den. Formlen er givet ved

\[
v = iR
\]

Med denne formel kan vi beregne spændingen \(v\), hvis værdierne af \(i\) og \(R\) er kendte. Hvis for eksempel \(i = 13 \, \text{A}\) og \(R = 5 \, \Omega\), så er

\[
v = iR = (13)(5) = 65
\]

Spændingen er altså 65 V.

Eksemplet viser, hvor vigtigt det er at holde nøje øje med enhederne for de fysiske størrelser, der indgår. En formel er kun gyldig, hvis man bruger et konsistent sæt enheder.

\begin{example} Indsættelse i formler \newline
Den kinetiske energi \(K\) for et legeme med massen \(M\), som bevæger sig med farten \(v\), kan beregnes med formlen:

\[
K = \frac{1}{2}Mv^2
\]

Beregn den kinetiske energi for et legeme med en masse på 5 kg, som bevæger sig med en fart på \(2 \, \text{m} \, \text{s}^{-1}\).

\begin{solution}
\[
K = \frac{1}{2}Mv^2 = \frac{1}{2}(5)(2^2) = 10
\] 

\end{solution}

I SI-systemet er enheden for energi joule, så legemets kinetiske energi er 10 joule.
\end{example}

\begin{example} Indsættelse i formler \newline
Arealet \(A\) af en cirkel med radius \(r\) kan beregnes med formlen \(A = \pi r^2\).

Hvis cirklens diameter \(d\) er kendt, kan man i stedet bruge den tilsvarende formel:

\[
A = \frac{\pi d^2}{4}
\]

Beregn arealet af en cirkel med en diameter på 0.1 m. Værdien af \(\pi\) er indbygget i din lommeregner.

\begin{solution}
    \[
A = \frac{\pi (0.1)^2}{4} = 0.00785 \, \text{m}^2
\]

\end{solution}


\end{example}

\begin{example} Indsættelse i formler \newline
Rumfanget \(V\) af en cirkulær cylinder er lig med dens tværsnitsareal \(A\) ganget med dens længde \(h\).

Beregn rumfanget af en cylinder med en diameter på 0.1 m og en længde på 0.3 m.

\begin{solution}
  \[
V = Ah = \frac{\pi (0.1)^2}{4} \times 0.3 = 0.00236
\]

Rumfanget er \(0.00236 \, \text{m}^3\).
\end{solution}
\end{example}

\section{Omskrivning af formler}
I formlen for cirklens areal, \(A = \pi r^2\), er variablen \(A\) isoleret. En variabel er isoleret, hvis den står alene på den ene side af lighedstegnet, normalt på venstre side, og ikke optræder andre steder i formlen. Hvis vi bliver bedt om at isolere \(r\) i formlen eller løse den med hensyn til \(r\), skal vi omskrive ligningen, så \(r\) bliver isoleret. Når man omskriver en formel, skal enhver operation, der udføres på den ene side, også udføres på den anden side. Der er fem centrale regler, som skal følges i denne proces.

\begin{custombox}{Regler for omskrivning af formler}
Følgende operationer kan udføres på begge sider af formlen:
\begin{itemize}
    \item Læg den samme størrelse til på begge sider
    \item Træk den samme størrelse fra på begge sider
    \item Gang begge sider med den samme størrelse forskellig fra nul -- husk at gange alle led
    \item Divider begge sider med den samme størrelse forskellig fra nul -- husk at dividere alle led
    \item Anvend den samme funktion på begge sider, og kontrollér samtidig funktionens definitionsmængde, og om operationen kan vendes om
\end{itemize}
\end{custombox}

De første fire operationer bevarer ækvivalensen, når deres angivne begrænsninger er opfyldt. At anvende en funktion, for eksempel at kvadrere, kan ikke nødvendigvis vendes om: For eksempel medfører \(x=2\), at \(x^2=4\), men \(x^2=4\) tillader også \(x=-2\). Løsninger, der er fundet efter en operation, som ikke kan vendes om, skal derfor kontrolleres i den oprindelige ligning.

\begin{example} Omskriv formlen \(p = 5t - 17\), så \(t\) bliver isoleret.

\begin{solution}
 For at isolere \(t\) på venstre side går vi frem trinvis efter de fem regler. Læg først 17 til på begge sider af ligningen \(p = 5t - 17\):
\[
p + 17 = 5t - 17 + 17
\]
Ved reduktion får vi:
\[
p + 17 = 5t
\]
Divider derefter begge sider med 5 for at isolere \(t\):
\[
\frac{p + 17}{5} = t
\]
Formlen for \(t\) er altså:
\[
t = \frac{p + 17}{5}
\]   
\end{solution}    
\end{example}

\begin{example} Omskriv formlen \(\sqrt{2q} = p\), og løs den med hensyn til \(q\).

\begin{solution}
   Kvadrér først begge sider for at fjerne kvadratroden omkring \(2q\). Bemærk, at \((\sqrt{2q})^2 = 2q\). Det giver:
\[
2q = p^2
\]
Divider derefter begge sider med 2 for at løse med hensyn til \(q\):
\[
q = \frac{p^2}{2}
\]   
\end{solution}
 
\end{example}

\newpage

\begin{example}
Omskriv formlen \(v = \sqrt{t^2 + w}\), og løs den med hensyn til \(w\). Følg disse trin for at isolere \(w\):   

\begin{solution}
    \begin{enumerate}[a.]
        \item Kvadrér først begge sider for at fjerne kvadratroden omkring \(t^2 + w\):
        \[
        v^2 = t^2 + w
        \]

        \item Træk derefter \(t^2\) fra på begge sider for at isolere \(w\):
        \[
        v^2 - t^2 = w
        \]

        \item Skriv til sidst formlen for \(w\):
        \[
        w = v^2 - t^2
        \]
    \end{enumerate}

    \end{solution}
\end{example}

\begin{example} Omskriv formlen \(x = \frac{1}{y}\), og løs den med hensyn til \(y\).

\begin{solution}
For at isolere \(y\) bemærker vi, at \(y\) står i nævneren. Hvis vi ganger begge sider med \(y\), forsvinder brøken:

\[
yx = y \times \frac{1}{y}
\]

Det reduceres til:

\[
yx = 1
\]

Divider til sidst begge sider med \(x\) for at løse med hensyn til \(y\):
\[
y = \frac{1}{x}
\]

Alternativt kan du blot invertere begge sider direkte og få:

\[
y = \frac{1}{x}
\]    
\end{solution}


\end{example}


\begin{example} Isolér \(R\) i formlen:

\[
\frac{2}{R} = \frac{3}{x + y}
\]

\begin{solution} Da \(R\) står i en brøk, inverterer vi begge sider:

\[
\frac{R}{2} = \frac{x + y}{3}
\]

Ved at gange begge sider med 2 får vi: $R = \dfrac{2(x + y)}{3}$    
\end{solution}
    
\end{example} 

\begin{example} Isolér \(R\) i formlen:

\[
\frac{1}{R} = \frac{1}{R_1} + \frac{1}{R_2}
\]
\begin{solution}
De to led på højre side kan samles:
\[
\frac{1}{R_1} + \frac{1}{R_2} = \frac{R_2 + R_1}{R_1R_2}
\]

Formlen bliver så:
\[
\frac{1}{R} = \frac{R_2 + R_1}{R_1R_2}
\]

Til sidst inverterer vi begge sider og får:
\[
R = \frac{R_1R_2}{R_2 + R_1}
\]

\end{solution}
\end{example}
